\chapter{Conclusion}
\label{ch:conclusion}

\section{Answers to the Four Research Questions}
\label{sec:9-1}

\paragraph{RQ1: GPU index construction and capacity.}
A complete JHQ index can be built and searched on one GPU without keeping
the raw database vectors on the device. Training and encoding run on the
GPU, with bounded batches and a two-pass build controlling temporary
storage. Stella's 17.78 million vectors require 67.81\,GiB in FP32;
JHQ-GPU's resident process occupies 20.44\,GiB, including context and
workspace. Index-cold construction takes 10.54--10.65\,s, with a sampled peak of
22.58\,GiB on the 32\,GB device (\cref{sec:7-7}).

\paragraph{RQ2: Representation-aware execution.}
JHQ's code structure supports small lookup tables, contiguous reads and
parallel candidate selection. Splitting each 256-entry primary table
into two 16-entry tables preserves the real-arithmetic score and is the
fastest tested table form in 22 of 24 settings. Packing codes into words
makes adjacent scanning threads read consecutive memory, while threads
can cooperate to refine one candidate. For small batches, splitting a
query across blocks adds useful parallelism: scheduling gives
$2.727\times$ throughput at batch 8 relative to the single-block baseline
with cooperative refinement. Grouped selection matches a CPU sort of the
GPU-computed scores in all 55,296 checked query instances; single-block
differences occur only at tied cut-offs
(\cref{sec:7-3,sec:7-6,sec:7-strict-selection}).

\paragraph{RQ3: Workload-level refinement calibration.}
Unlabelled queries can select a reusable refinement budget that meets
the tested quality target, but sample size matters. Using all 500
calibration queries keeps held-out recall within $10^{-3}$ of the
highest measured recall across tested budgets in all twelve settings;
the maximum observed loss is 0.0006. Smaller samples can miss queries
that need more refinement: with 128 queries, arXiv exceeds the target
in 40 of 64 calibration samples at 128 probes and 32 of 64 at 512 probes. With 128-query calibration, direct OpenAI tests recover overhead
within 4--9 batches relative to the fixed-budget controls. On arXiv, the selected budget instead
spends more time to reach the quality target (\cref{sec:7-5}).

\paragraph{RQ4: Cross-system search performance.}
The main comparison favours JHQ-GPU at high batch throughput. At recall
0.97 and batch 1,000, it delivers $1.32$--$1.75\times$ the throughput of
CAGRA fp32 across four datasets. Across the tested grid, net process memory excluding CUDA context is
0.29--0.39 of CAGRA's index allocation. It leads the tested methods on Vogue and
both OpenAI datasets, while arXiv is nearly tied with IVF-RaBitQ;
Vogue's lead depends on grouped scanning (\cref{sec:7-2}). The OpenAI
measurements use TF32 multiplication; FP32 checks retain the lead over
CAGRA fp32 using the same trained model
(\cref{sec:b-precision-anchors}). This advantage has a batch-size boundary:
at batch 8 and recall 0.97, CAGRA fp32 is faster on all four datasets
(\cref{sec:7-smallbatch}).

\section{Contributions, Limitations and Outlook}
\label{sec:9-2}

Compression delivers capacity and speed when codes are cheap to access,
parallelism matches the batch and refinement meets the workload's
quality needs.

The evidence covers one GPU, six prepared datasets, squared-L2 top-10
search and static indexes. Historical CPU comparisons do not quantify a
speedup for the evaluated version (\cref{sec:7-cpu}). Reference agreement
on the calibration pool leaves future-query quality uncertain.

Two extensions follow. Risk-controlled calibration should account for
workload drift and expected reuse, and be compared with a quality-matched
fixed budget on total time, including calibration. Fetching residual
codes from host memory could expand capacity; transfer volume, overlap
and tail latency would determine when it pays off (\cref{sec:8-4}).
